\begin{definition}[Translated edge windows]%
    \label{def:peps_normalSquareTranslatedEdgeWindow}
    \lean{TNLean.PEPS.NormalSquareTranslatedEdgeWindow}
    \lean{TNLean.PEPS.NormalSquareTranslatedEdgeWindow.blockingDatum}
    \lean{TNLean.PEPS.NormalSquareTranslatedEdgeWindow.injective_chain}
    \lean{TNLean.PEPS.NormalSquareTranslatedEdgeWindow.endpoint_disjoint_cover}
    \lean{TNLean.PEPS.normalSquareTranslatedEdgeWindow_rightEdge_margins}
    \lean{TNLean.PEPS.normalSquareTranslatedEdgeWindow_upEdge_margins}
    \lean{TNLean.PEPS.not_leftBoundaryRightEdge_window}
    \lean{TNLean.PEPS.not_leftBoundaryRightEdge_window_seven}
    \lean{TNLean.PEPS.not_rightMarginRightEdge_window}
    \lean{TNLean.PEPS.not_rightMarginRightEdge_window_seven}
    \lean{TNLean.PEPS.not_bottomBoundaryUpEdge_window}
    \lean{TNLean.PEPS.not_bottomBoundaryUpEdge_window_seven}
    \lean{TNLean.PEPS.not_topMarginUpEdge_window}
    \lean{TNLean.PEPS.not_topMarginUpEdge_window_seven}
    \lean{TNLean.PEPS.not_fourSidedBoundaryWindow_seven}
    \lean{TNLean.PEPS.not_forall_normalSquareTranslatedEdgeWindow_seven}
    \lean{TNLean.PEPS.normalSquareHorizontalTranslatedEdgeWindow}
    \lean{TNLean.PEPS.normalSquareVerticalTranslatedEdgeWindow}
    \lean{TNLean.PEPS.normalSquareHorizontalTranslatedEdge_sub_eq_rightEdge}
    \lean{TNLean.PEPS.normalSquareVerticalTranslatedEdge_sub_eq_upEdge}
    \lean{TNLean.PEPS.squareLatticeRightEdgeWindow}
    \lean{TNLean.PEPS.squareLatticeUpEdgeWindow}
    \lean{TNLean.PEPS.squareLatticeRightEdgeWindow_of_margins}
    \lean{TNLean.PEPS.squareLatticeUpEdgeWindow_of_margins}
    \lean{TNLean.PEPS.horizontalSquareLatticeEdgeWindow}
    \lean{TNLean.PEPS.horizontalSquareLatticeEdgeWindow_of_margins}
    \lean{TNLean.PEPS.verticalSquareLatticeEdgeWindow}
    \lean{TNLean.PEPS.verticalSquareLatticeEdgeWindow_of_margins}
    \lean{TNLean.PEPS.IsNormalSquareHorizontalEdgeMargins}
    \lean{TNLean.PEPS.IsNormalSquareVerticalEdgeMargins}
    \lean{TNLean.PEPS.NormalSquareEdgeMarginCover}
    \lean{TNLean.PEPS.NormalSquareEdgeMarginCover.window}
    \lean{TNLean.PEPS.NormalSquareEdgeMarginCover.blockingDatum}
    \lean{TNLean.PEPS.NormalSquareEdgeMarginCover.injective_chain}
    \lean{TNLean.PEPS.NormalSquareEdgeMarginCover.endpoint_disjoint_cover}
    \lean{TNLean.PEPS.normalSquareEdgeMarginCover_rightEdge_bounds}
    \lean{TNLean.PEPS.normalSquareEdgeMarginCover_upEdge_bounds}
    \lean{TNLean.PEPS.not_leftBoundaryRightEdge_marginCover}
    \lean{TNLean.PEPS.not_leftBoundaryRightEdge_marginCover_seven}
    \lean{TNLean.PEPS.not_rightMarginRightEdge_marginCover}
    \lean{TNLean.PEPS.not_rightMarginRightEdge_marginCover_seven}
    \lean{TNLean.PEPS.not_bottomBoundaryUpEdge_marginCover}
    \lean{TNLean.PEPS.not_bottomBoundaryUpEdge_marginCover_seven}
    \lean{TNLean.PEPS.not_topMarginUpEdge_marginCover}
    \lean{TNLean.PEPS.not_topMarginUpEdge_marginCover_seven}
    \lean{TNLean.PEPS.not_fourSidedBoundaryMarginCover_seven}
    \leanok
    \uses{lem:peps_normalSquareHorizontalTranslatedEdge_blockingData,
        lem:peps_normalSquareVerticalTranslatedEdge_blockingData}
    A translated edge window records that a chosen edge is one of the
    translated horizontal or vertical edge pictures, together with a
    rectangular cover of the complementary block. Such a window gives the
    one-edge blocking datum for that edge, hence the three injective regions
    used in the edge-blocked chain. The coordinate right and upward edges
    obtain such windows directly from the corresponding translated picture.
    Equivalently, an actual coordinate right or upward edge obtains such a
    window when it has the necessary room around it for the translated
    $5\times7$ or $7\times5$ frame and the complementary cover is supplied.
    The same criterion applies to an arbitrary horizontal or vertical
    square-lattice edge after identifying it with its coordinate right or
    upward representative. The oriented margin-and-cover datum is the same
    criterion written as per-edge data in the rectangular coordinate model;
    it directly supplies the one-edge blocking datum, its three injective
    regions, and the identities
    $u_e\in R_e$, $v_e\in B_e$,
    $R_e\cap B_e=R_e\cap C_e=B_e\cap C_e=\varnothing$, and
    $R_e\cup B_e\cup C_e=V$. The boundary obstruction lemmas show that these
    translated-window and margin-cover criteria are interior criteria, not the
    final every-edge construction in the open square-lattice model.
    % Detailed source-comparison status:
    % docs/paper-gaps/peps_normal_ft_section3_route.tex.
\end{definition}

\begin{theorem}[Construction from translated edge windows]%
    \label{thm:peps_normalSquareTranslatedEdgeBlockingHypotheses_of_windows}
    \lean{TNLean.PEPS.normalSquareTranslatedEdgeBlockingHypotheses_of_windows}
    \lean{TNLean.PEPS.normalSquareTranslatedEdgeBlockingHypotheses_blockingData_of_windows}
    \lean{TNLean.PEPS.normalSquareTranslatedEdgeBlockingHypotheses_injective_chain_of_windows}
    \lean{TNLean.PEPS.%
        normalSquareTranslatedEdgeBlockingHypotheses_endpoint_disjoint_cover_of_windows}
    \lean{TNLean.PEPS.normalSquareEdgeBlockingHypotheses_of_marginCovers}
    \lean{TNLean.PEPS.normalSquareEdgeBlockingHypotheses_blockingData_of_marginCovers}
    \lean{TNLean.PEPS.normalSquareEdgeBlockingHypotheses_injective_chain_of_marginCovers}
    \lean{TNLean.PEPS.normalSquareEdgeBlockingHypotheses_endpoint_disjoint_cover_of_marginCovers}
    \lean{TNLean.PEPS.not_forall_normalSquareEdgeMarginCover_seven}
    \leanok
    \uses{def:peps_normalSquareTranslatedEdgeWindow,
        def:peps_normalEdgeBlockingHypotheses}
    Given rectangular injectivity hypotheses, finite union closure, and a
    translated edge window at every edge of the finite square lattice, the
    corresponding one-edge blockings determine the normal edge-blocking
    hypotheses. This is the conditional construction;
    constructing such windows for all edges of the $7\times7$ lattice remains
    separate; the current open rectangular translated-window criterion does not
    by itself supply such a family. The one-edge datum recovered from the
    resulting hypotheses is the datum supplied by the corresponding translated
    window, and the resulting hypotheses give the three injective regions at
    every edge, together with
    $u_e\in R_e$, $v_e\in B_e$,
    $R_e\cap B_e=R_e\cap C_e=B_e\cap C_e=\varnothing$, and
    $R_e\cup B_e\cup C_e=V$. Equivalently, it is enough to supply the oriented
    margin-and-cover data at every edge; the resulting hypotheses recover the
    corresponding datum, give the same three injectivity conclusions, and give
    the same endpoint, disjointness, and covering conclusions. The obstruction
    theorem records that this interior margin-cover criterion does not by
    itself supply data at every edge of the open $7\times7$ lattice.
    % Detailed source-comparison status:
    % docs/paper-gaps/peps_normal_ft_section3_route.tex.
\end{theorem}
\begin{proof}\leanok
    Apply the general construction for one-edge blocking data to the datum
    supplied by the chosen translated window at each edge.
\end{proof}

\begin{lemma}[Conditional injectivity of the normal $R$ and $S$ regions]%
    \label{lem:peps_normalSquareRegionRS_injective_of_unionClosure}
    \lean{TNLean.PEPS.NormalSquareLatticeRectangleInjectivityHypotheses.regionR_injective_of_union}
    \lean{TNLean.PEPS.NormalSquareLatticeRectangleInjectivityHypotheses.regionS_injective_of_union}
    \leanok
    \uses{def:peps_regionInjectivityData,
        def:peps_normalRectangleInjectivityHypotheses,
        lem:peps_normalSquareRegionRS_rectangular_decomposition}
    If region injectivity is closed under unions, then the coordinate
    square-lattice rectangular injectivity hypotheses imply that the displayed
    regions $R$ and $S$ are injective.
\end{lemma}
\begin{proof}\leanok
    Region $R$ is the union of one injective $2\times3$ rectangle and one
    injective $3\times2$ rectangle. Region $S$ is the union of two
    injective $2\times3$ rectangles. Apply union closure in each case.
\end{proof}

\begin{lemma}[Injectivity of the normal $R$, $S$, and $T$ regions]%
    \label{lem:peps_normalSquareRegionRST_injective}
    \lean{TNLean.PEPS.regionBlockedTensorInjective_normalSquareRegionR}
    \lean{TNLean.PEPS.regionBlockedTensorInjective_normalSquareRegionS}
    \lean{TNLean.PEPS.regionBlockedTensorInjective_normalSquareRegionT}
    \leanok
    \uses{thm:peps_regionInjectivityUnionClosure_of_overlap,
        lem:peps_normalSquareRegionRS_injective_of_unionClosure,
        lem:peps_normalSquareRegionT_injective_of_cover}
    Fix a square-lattice PEPS tensor with every bond dimension positive whose
    blocked tensor is injective on every contiguous $2\times3$ and $3\times2$
    rectangle. Then the displayed regions $R$ and $S$ are injective. If
    moreover the displayed region $T$ has a cover by such rectangles, then $T$
    is injective too. The union closure used here is supplied by the
    overlapping union lemma, so no single-vertex injectivity is assumed.
\end{lemma}
\begin{proof}\leanok
    Apply the conditional injectivity of $R$ and $S$
    (Lemma~\ref{lem:peps_normalSquareRegionRS_injective_of_unionClosure}) and the
    conditional cover criterion for $T$
    (Lemma~\ref{lem:peps_normalSquareRegionT_injective_of_cover}), with the union
    closure discharged by
    Theorem~\ref{thm:peps_regionInjectivityUnionClosure_of_overlap}.
\end{proof}

% The comparison square and comparison region at a parametric offset, used to
% place the one-site comparison pair $R\subseteq S=R\cup\{v\}$ afresh at every
% vertex of the open square lattice.

\begin{definition}[Comparison square, comparison region, and interior window]%
    \label{def:peps_normalSquareComparisonRegions}
    \lean{TNLean.PEPS.normalSquareComparisonSquare}
    \lean{TNLean.PEPS.normalSquareComparisonRegion}
    \lean{TNLean.PEPS.IsNormalSquareComparisonWindow}
    \leanok
    \uses{def:peps_squareLatticeCoordinateRegions}
    At a coordinate offset $(a,b)$ the comparison square is the $3\times3$
    contiguous coordinate rectangle with lower-left corner $(a,b)$. The
    comparison region is this square with its lower-left corner $(a,b)$
    removed, presented as the union of the $2\times3$ contiguous rectangle at
    $(a+1,b)$ and the $3\times2$ contiguous rectangle at $(a,b+1)$. This is an
    auxiliary one-site comparison pair carrying its omitted site at the
    lower-left corner; it is the reflected counterpart of the displayed regions
    $R$ and $S$, whose omitted site $S\setminus R$ sits at the lower-right
    corner, placed at an arbitrary offset rather than at a fixed origin. The
    interior window at $(a,b)$ for a finite $n\times m$ lattice is the placement
    margin
    $4\le a$, $a+9\le n$, $4\le b$, $b+9\le m$, under which every boundary edge
    of the comparison square and of the comparison region carries the interior
    margins of the open edge-blocking geometry.
\end{definition}

\begin{lemma}[Membership in the comparison region]%
    \label{lem:peps_normalSquareComparisonRegion_membership}
    \lean{TNLean.PEPS.mem_normalSquareComparisonRegion}
    \lean{TNLean.PEPS.notMem_normalSquareComparisonRegion}
    \lean{TNLean.PEPS.insert_normalSquareComparisonRegion}
    \leanok
    \uses{def:peps_normalSquareComparisonRegions,
        lem:peps_squareLatticeCoordinateRegion_identities}
    For any offset $(a,b)$ a vertex with coordinates $(x,y)$ lies in the
    comparison region at $(a,b)$ if and only if
    \begin{align}
        &(a+1\le x<a+3\wedge b\le y<b+3)
        \lor
        (a\le x<a+3\wedge b+1\le y<b+3).
        \label{eq:peps_comp_region_mem}
    \end{align}
    Let $v$ be a lattice vertex and place the comparison region at $v$'s own
    coordinates. Then $v$ lies outside that region, and inserting $v$ back in
    recovers the full $3\times3$ comparison square placed at $v$'s coordinates.
\end{lemma}
\begin{proof}\leanok
    Membership unfolds the union of the two contiguous rectangles into
    (\ref{eq:peps_comp_region_mem}). Place the region at $v$'s own
    coordinates, so $(a,b)=(x,y)=(v_1,v_2)$. The corner offset $(v_1,v_2)$ fails
    the first disjunct because $v_1+1\le v_1$ is false, and fails the second
    because $v_2+1\le v_2$ is false:
    \begin{align}
        &\neg\bigl((v_1+1\le v_1<v_1+3)
            \wedge(v_2\le v_2<v_2+3)\bigr)
        \notag\\
        &\qquad\wedge
        \neg\bigl((v_1\le v_1<v_1+3)
            \wedge(v_2+1\le v_2<v_2+3)\bigr).
        \notag
    \end{align}
    Thus $v$ lies outside the region. Inserting $v$ therefore yields the region
    together with $\{(v_1,v_2)\}$, and a vertex lies in this set exactly when it
    equals $v$ or already lies in the region. That is the membership condition
    of the $3\times3$ square, the corner being the one cell of the square that
    both pieces omit.
\end{proof}

\begin{lemma}[Band decomposition of the comparison square complement]%
    \label{lem:peps_normalSquareComparisonSquare_complementBands}
    \lean{TNLean.PEPS.compl_normalSquareComparisonSquare_eq_bands}
    \leanok
    \uses{def:peps_normalSquareComparisonRegions}
    On a finite $n\times m$ lattice, if $a+3\le n$ and $b+3\le m$ then the
    complement of the comparison square at $(a,b)$ is the union of the four
    coordinate bands
    \begin{align}
        &\{(x,y):0\le x<a,\ 0\le y<m\}
        \cup
        \{(x,y):a+3\le x<n,\ 0\le y<m\}
        \notag\\
        &\qquad\cup
        \{(x,y):a\le x<a+3,\ 0\le y<b\}
        \cup
        \{(x,y):a\le x<a+3,\ b+3\le y<m\}.
        \label{eq:peps_comp_square_bands}
    \end{align}
\end{lemma}
\begin{proof}\leanok
    A vertex $(x,y)$ lies in the comparison square exactly when
    $a\le x<a+3$ and $b\le y<b+3$. With $0\le x<n$ and $0\le y<m$ on the
    lattice, the negation of that conjunction is membership in
    (\ref{eq:peps_comp_square_bands}):
    \begin{align}
        &\neg(a\le x<a+3\wedge b\le y<b+3)
        \notag\\
        &\quad\Longleftrightarrow
        (0\le x<a\wedge 0\le y<m)
        \lor
        (a+3\le x<n\wedge 0\le y<m)
        \notag\\
        &\qquad\lor
        (a\le x<a+3\wedge 0\le y<b)
        \lor
        (a\le x<a+3\wedge b+3\le y<m).
        \notag
    \end{align}
\end{proof}

\begin{lemma}[Band decomposition of the comparison region complement]%
    \label{lem:peps_normalSquareComparisonRegion_complementBands}
    \lean{TNLean.PEPS.compl_normalSquareComparisonRegion_eq_bands}
    \leanok
    \uses{def:peps_normalSquareComparisonRegions,
        lem:peps_normalSquareComparisonRegion_membership}
    On a finite $n\times m$ lattice, if $2\le a$, $1\le b$, $a+3\le n$, and
    $b+3\le m$ then the complement of the comparison region at $(a,b)$ is the
    union of the same four bands of
    Lemma~\ref{lem:peps_normalSquareComparisonSquare_complementBands} together
    with the $3\times2$ corner rectangle
    $\{(x,y):a-2\le x<a+1,\ b-1\le y<b+1\}$,
    which contains the omitted corner $(a,b)$.
\end{lemma}
\begin{proof}\leanok
    By (\ref{eq:peps_comp_region_mem}), a vertex $(x,y)$ lies in the complement
    exactly when both disjuncts fail:
    \begin{align}
        &\neg(a+1\le x<a+3\wedge b\le y<b+3)
        \notag\\
        &\qquad\wedge
        \neg(a\le x<a+3\wedge b+1\le y<b+3).
        \notag
    \end{align}
    Together with $0\le x<n$ and $0\le y<m$, this places the vertex either in
    one of the four bands of
    (\ref{eq:peps_comp_square_bands}) or in the $3\times2$ corner rectangle.
    Writing $R$ for the comparison region and $S$ for the comparison square,
    \begin{align}
        V\setminus R
        &=
        \{(x,y):a-2\le x<a+1,\ b-1\le y<b+1\}
        \cup(V\setminus S),
        \notag
    \end{align}
    since the comparison region is the comparison square with the corner $(a,b)$
    removed, and that corner lies in the displayed $3\times2$ rectangle. The
    hypotheses $2\le a$ and $1\le b$ keep the corner rectangle's lower-left
    coordinates $(a-2,b-1)$ within the lattice.
\end{proof}

\begin{lemma}[Injectivity of the comparison region and the comparison square]%
    \label{lem:peps_normalSquareComparison_injective}
    \lean{TNLean.PEPS.NormalSquareLatticeRectangleInjectivityHypotheses.comparisonRegion_injective}
    \lean{TNLean.PEPS.NormalSquareLatticeRectangleInjectivityHypotheses.comparisonSquare_injective}
    \leanok
    \uses{def:peps_normalRectangleInjectivityHypotheses,
        def:peps_regionInjectivityData,
        lem:peps_squareLatticeRectangleInjectivity_rectangles,
        lem:peps_normalSquareWideRectangle_injective}
    Assume rectangular injectivity and union closure. On a finite $n\times m$
    lattice, if $a+3\le n$ and $b+3\le m$ then the comparison region at $(a,b)$
    is injective, being the union of a $2\times3$ and a $3\times2$ rectangle,
    and the comparison square at $(a,b)$ is injective, being a $3\times3$
    rectangle.
\end{lemma}
\begin{proof}\leanok
    Apply union closure to the rectangular injectivity of the contained
    $2\times3$ and $3\times2$ rectangles for the region, the square being a
    single large rectangle.
\end{proof}

\begin{lemma}[Injectivity of the comparison complements]%
    \label{lem:peps_normalSquareComparison_complement_injective}
    \lean{TNLean.PEPS.NormalSquareLatticeRectangleInjectivityHypotheses.compl_comparisonSquare_injective}
    \lean{TNLean.PEPS.NormalSquareLatticeRectangleInjectivityHypotheses.compl_comparisonRegion_injective}
    \leanok
    \uses{def:peps_normalRectangleInjectivityHypotheses,
        def:peps_regionInjectivityData,
        lem:peps_normalSquareWideRectangle_injective,
        lem:peps_normalSquareComparisonSquare_complementBands,
        lem:peps_normalSquareComparisonRegion_complementBands}
    Assume rectangular injectivity and union closure. On a finite $n\times m$
    lattice, if $2\le a$, $2\le b$, $a+5\le n$, and $b+5\le m$ then the
    complement of the comparison square at $(a,b)$ is injective, being the union
    of its four surrounding bands, and the complement of the comparison region
    at $(a,b)$ is injective, being that union together with one further corner
    rectangle.
\end{lemma}
\begin{proof}\leanok
    Rewrite each complement by its band decomposition and apply union closure to
    the bands, each a large contiguous rectangle, and to the corner rectangle.
\end{proof}

\begin{definition}[Normal edge-blocking hypotheses]%
    \label{def:peps_normalEdgeBlockingHypotheses}
    \lean{TNLean.PEPS.NormalEdgeBlockingData}
    \lean{TNLean.PEPS.NormalEdgeBlockingHypotheses}
    \leanok
    \uses{def:peps_edge, def:peps_regionInjectivityData}
    Around one edge, the normal proof uses a blocking into three injective
    regions: a red region containing one endpoint, a blue region containing the
    other endpoint, and an injective complementary region. The three regions
    are pairwise disjoint and cover the whole vertex set. The edge-blocking
    hypotheses are precisely this one-edge datum at every edge; the red,
    blue, and complementary regions are read off from it. The local blocking
    picture used in the proof of
    \cite[Theorem~3]{Molnar2018NormalPEPS}:
    \begin{tenkzequation}
        % Formula verdict: in the source 7-by-7 frame, A1=(4-6,2-3),
        % A2=(3-4,4-6), and A3 is their cell-set complement; the marked
        % horizontal edge is (4,3)-(4,4), joining one endpoint in each block.
        % Ink-to-index verdict: every dot is an ordinary tensor, every lattice
        % edge is a virtual contraction, the three fills encode the disjoint
        % cover, and the distinguished edge is the edge being blocked.
        % Contraction/boundary verdict: the panel has 84 internal virtual
        % contractions and 28 cut virtual bonds, with no exposed physical
        % indices, hence boundary type (V,P)=(28,0).
        % Source verdict: exact discrete translation of arXiv:1804.04964,
        % paper_normal.tex:1475--1499; A1 is red, A2 blue, and A3 the rest.
        \begin{tenkzlattice}[rows=7, cols=7, boundary legs]
            \tnregion[slot=selected, label={$A_1$},
                label at=center]{(4-6,2-3)}
            \tnregion[slot=secondary, label={$A_2$},
                label at=center]{(3-4,4-6)}
            \tnregion[slot=complement, outline, label={$A_3$},
                label at=north]
                {(1-7,1-7) - (4-6,2-3) - (3-4,4-6)}
            \tnedge[distinguished]{(4,3)-(4,4)}
        \end{tenkzlattice}
    \end{tenkzequation}
\end{definition}

\begin{theorem}[Injectivity supplied by normal edge blocking]%
    \label{thm:peps_normalEdgeBlockingHypotheses_injectiveChain}
    \lean{TNLean.PEPS.NormalEdgeBlockingData.injective_chain}
    \lean{TNLean.PEPS.NormalEdgeBlockingData.endpoint_disjoint_cover}
    \lean{TNLean.PEPS.NormalEdgeBlockingData.endpoint_injective_disjoint_cover}
    \lean{TNLean.PEPS.NormalEdgeBlockingHypotheses.ofBlockingData}
    \lean{TNLean.PEPS.NormalEdgeBlockingHypotheses.ofBlockingData_blockingData}
    \lean{TNLean.PEPS.NormalEdgeBlockingHypotheses.injective_chain_at_edge}
    \lean{TNLean.PEPS.NormalEdgeBlockingHypotheses.endpoint_disjoint_cover_at_edge}
    \lean{TNLean.PEPS.%
        NormalEdgeBlockingHypotheses.endpoint_injective_disjoint_cover_at_edge}
    \leanok
    \uses{def:peps_normalEdgeBlockingHypotheses}
    A family of one-edge blocking data is precisely the normal edge-blocking
    hypotheses. If $H$ is assembled from the family $d_e$, then the one-edge
    datum that $H$ associates with $e$ is $d_e$. Write the endpoints of $e$
    as $u_e<v_e$ in the canonical vertex ordering. If the three regions are
    $R_e$, $B_e$, and $C_e$, then
    \begin{align}
        u_e&\in R_e,
        \notag\\
        v_e&\in B_e,
        \notag\\
        R_e\cap B_e&=R_e\cap C_e=B_e\cap C_e=\varnothing,
        \notag\\
        R_e\cup B_e\cup C_e&=V,
        \label{eq:peps_edge_blocking_cover}
    \end{align}
    and the underlying PEPS tensor is injective on $R_e$, $B_e$, and $C_e$.
\end{theorem}
\begin{proof}\leanok
    By construction, the one-edge datum at each edge $e$ is $d_e$.
    The endpoint, disjointness, and coverage statements in
    (\ref{eq:peps_edge_blocking_cover}), together with the injectivity
    statements, are exactly the corresponding assumptions of that one-edge
    datum.
\end{proof}
